\begin{frame}[t]{Graph measures link protein binding to actomyosin self-organization.}
\vspace{0pt}
\begin{center}\begin{tikzpicture}[x=1cm,y=1cm,every node/.style={align=center,font=\fontsize{12}{14.5}\selectfont}]
\node[state,text width=3.9cm,minimum height=1.55cm] (sim) at (0,0) {\textbf{Cytosim}\\Driven active matter\\Filaments + motors};
\node[state,text width=3.9cm,minimum height=1.55cm] (graph) at (4.8,0) {\textbf{Connectivity graph}\\Nodes: filaments\\Edges: protein links};
\node[evidence,text width=3.9cm,minimum height=1.55cm] (result) at (9.6,0) {\textbf{Graph readout}\\Cliques; communities\\Bundles and clusters};
\draw[flow](sim)--(graph);\draw[flow](graph)--(result);
\node[text width=13.3cm,font=\fontsize{12}{15}\selectfont] at (4.8,-1.65) {\textbf{2--7} binding sites per linker\quad \textbf{600 s} simulated time\quad \textbf{30} random starts};
\node[text width=13.3cm,text=Teal,font=\fontsize{13}{16}\selectfont] at (4.8,-2.55) {More binding sites promoted bundles and larger clusters.};
\end{tikzpicture}\end{center}
\tagline{Biological motivation: contraction and dendritic-spine plasticity involved in memory.}
\sources{Eliaz, Nedelec, Morrison, Levine \& Cheung, \href{https://doi.org/10.1103/PhysRevE.102.062420}{Phys. Rev. E 102:062420 (2020)}.}
\qadetail{Biological context is distinct from a simulated functional outcome. The paper motivates actin/CaMKII architecture through neuronal dendritic-spine plasticity and considers active contraction from ATP-powered motors; it does not simulate a memory task or a whole muscle. Graph theory provides local and global order parameters rather than treating every connected cluster as a homogeneous morphology. Protein-mediated filament connections define the graph; clique, modularity/community and connectivity measures answer different structural questions. The footer omits the requested authorship commentary. Primary author preprint: https://arxiv.org/abs/2006.06503, Insights from graph theory on the morphologies of actomyosin networks with multilinkers, Yossi Eliaz, Francois Nedelec, Greg Morrison, Herbert Levine and Margaret S. Cheung. The study extends Cytosim's stochastic semiflexible-filament model with multivalent actin-binding proteins. Filaments become graph nodes and protein-mediated filament connections become edges. Linker valencies are 2,3,4,5,6,7. The reported low/high motor-content comparisons use 600-second simulations and statistics from 30 simulations with different random initial conditions. Do not infer 30 replicates for every condition in the entire parameter sweep. Higher valency promotes bundles and large clusters under the studied conditions; active myosin motors also generate dendritic branches from bundles. This is a physical-network simulation finding, not a fork-speed or agent-dependence result. The user's confirmed research path begins in biological physics. The advisor's group record confirms a Physics PhD at the University of Houston in September 2020: https://sites.google.com/nsm.uh.edu/cheung/past-group-members. The supplied earlier deck records association with Rice's Center for Theoretical Biological Physics and the thesis Physics of self-assembly in complex matter (2020), later self-published as Physics of Complex (Biological) Matter (2023); the exact thesis title was not independently located in a university repository. Related coauthored work: Liman et al., The role of the Arp2/3 complex in shaping the dynamics and structures of branched actomyosin networks, PNAS 117:10825 (2020); Li, Liman, Eliaz and Cheung, Forecasting avalanches in branched actomyosin networks with network science and machine learning, J. Phys. Chem. B 125:11591--11605 (2021), https://arxiv.org/abs/2106.04005 and DOI 10.1021/acs.jpcb.1c04792. The latter varies Arp2/3 concentration in mechanochemical simulations: high concentration stalls networks, low concentration contracts them, and intermediate concentration permits loosely connected clusters to collapse suddenly. Graph features and machine learning characterize and forecast these changes. The earlier bibliography also retains the coauthored protein-folding-temperature study, Zegarra et al., Phys. Rev. E 97:032402 (2018), and CaXML, Zhang, Nde, Eliaz, Jennings, Cieplak and Cheung, Protein Science 34:e70023 (2025). These works establish methodological background, not measured sandbox behavior.}
\note{[0:50] Actomyosin networks are active matter: myosin motors consume ATP and pull on actin filaments. These networks support contraction, while actin organization in neuronal spines is relevant to memory. We extended Cytosim with proteins that bind two to seven filaments. A graph represents each filament as a node and protein connections as edges. Clique and community measures distinguish local bundles from larger connected structures. In the reported comparison, we simulated six hundred seconds with thirty random starts. More binding sites promoted bundles and larger clusters. The graph made self-organization measurable; genomics then made its representation a computational bottleneck.}
\end{frame}
